% Part: computability
% Chapter: recursive-functions
% Section: primitive-recursion

\documentclass[../../../include/open-logic-section]{subfiles}

\begin{document}

% SEGMENT OLP-0211-S01

\olfileid{cmp}{rec}{pre}
\olsection{முதன்மை மீள்வரையறை}

அடுத்தெண் செயல்~$+1$ ஐ இறுதிக்குட்பட்ட எண்ணிக்கையில் பயன்படுத்தி
ஒவ்வொரு இயல் எண்ணையும்~$0$ இலிருந்து அடையலாம் என்பது இயல் எண்களின்
ஒரு சிறப்பியல்பு---எந்த இயல் எண்ணும்~$0$ அல்லது $0$ இன் அடுத்தெண்
\dots{} அடுத்தெண். இவ்வுண்மையைப் பயன்படுத்தும்
சார்பு~$h\colon\Nat \to \Nat$ ஒன்றைக் குறிப்பிடுவதற்கான ஒரு வழி:
(a)~உள்ளீடு~$0$ க்கான $h$ இன் மதிப்பு என்னவென்று குறிப்பிடுக;
(b)~$h(x)$ இன் மதிப்பு கொடுக்கப்பட்டால், $h(x+1)$ இன் மதிப்பை
எவ்வாறு கணிப்பது என்றும் குறிப்பிடுக. ஏனெனில் (a), $h(0)$ என்னவென்று
நேரடியாகக் கூறுவதால், $h$ ஆனது~$0$ க்கு வரையறுக்கப்பட்டுள்ளது.
இப்போது (b) தந்த அறிவுறுத்தலை $x=0$ க்குப் பயன்படுத்தி,
~$h(0)$ இலிருந்து $h(1) = h(0+1)$ ஐக் கணிக்கலாம். அதே
அறிவுறுத்தலை $x=1$ க்குப் பயன்படுத்தி,~$h(1)$ இலிருந்து
$h(2) = h(1+1)$ ஐக் கணிக்கிறோம்; இவ்வாறே தொடர்கிறது. ஒவ்வொரு
இயல் எண்~$x$ க்கும், $h(x+1)$ இலிருந்து $h(x)$ ஐ வரையறுக்கும்
படியை இறுதியில் அடைவோம்; எனவே அனைத்து $x \in \Nat$ க்கும்
$h(x)$ வரையறுக்கப்பட்டுள்ளது.

எடுத்துக்காட்டாக, பின்வரும் இரு சமன்பாடுகளால்
$h\colon \Nat \to \Nat$ ஐக் குறிப்பிடுகிறோம் என்று கொள்க:
\begin{align*}
h(0) & =  1\\
h(x+1) & =  2 \cdot h(x)
\end{align*}
எவ்வாறு பெருக்குவது என்பது ஏற்கெனவே தெரிந்தால், மேலுள்ள (a) மற்றும்
~(b) க்குத் தேவையான தகவல்களை இச்சமன்பாடுகள் தருகின்றன. இரண்டாம்
சமன்பாட்டை அடுத்தடுத்து பயன்படுத்துவதால் பின்வருவன கிடைக்கின்றன:
\begin{align*}
  h(1) & = 2\cdot h(0) = 2,\\
  h(2) & = 2\cdot h(1) = 2\cdot 2,\\
  h(3) & = 2 \cdot h(2) = 2\cdot 2 \cdot 2,\\
  & \vdots
\end{align*}
நாம் குறிப்பிட்ட சார்பு~$h$ ஆனது $h(x) = 2^x$ என்பதைக் காண்கிறோம்.

இவ்விரு நெறிகளையும் நிறைவுறுத்தும் ஒரே சார்பு~$h$ மட்டுமே உள்ளது
என்பதை இயல் எண்களின் சிறப்பியல்பு உறுதிசெய்கிறது. இவற்றைப் போன்ற
ஒரு சோடி சமன்பாடுகள், சார்பு~$h$ இன்
\emph{முதன்மை மீள்வரையறை வழி வரையறை} எனப்படும். நாம்~$h$ ஐ
“மீள்வரையறை முறையில்” வரையறுப்பதால் இப்பெயர் பெறுகிறது; அதாவது,
வரையறையில், குறிப்பாக இரண்டாம் சமன்பாட்டின் வலப்பக்கத்தில்,
$h$ தானே இடம்பெறுகிறது. $h(x+1)$ ஐ வரையறுக்கும்போது அதற்கு
உடனடியாக முந்திய மதிப்பு~$h(x)$ ஐ மட்டுமே பயன்படுத்துவதால் அது
“முதன்மை” ஆகும். ~$\Nat$ மீது ஒரு சார்பை மீள்வரையறை முறையில்
வரையறுப்பதற்கான மிக எளிய வழி இது.

கூட்டல், பெருக்கல் போன்ற இன்னும் அடிப்படையான சார்புகளையும் முதன்மை
மீள்வரையறையால் வரையறுக்கலாம். எனினும் இந்நேர்வுகளில் தொடர்புடைய
சார்புகள் $2$-இடச் சார்புகள். உள்ளீட்டு இடங்களில் ஒன்றை நிலையாகக்
கொண்டு, மற்றொன்றை மீள்வரையறைக்குப் பயன்படுத்துகிறோம். எடுத்துக்காட்டாக,
$\Add(x, y)$ ஐ வரையறுக்க,~$x$ ஐ நிலையாகக் கொண்டு, முதலில்
$y=0$ க்கான மதிப்பையும், பின்னர்~$y$ ஐக் கொண்டு $y+1$ க்கான
மதிப்பையும் வரையறுக்கலாம். $x$ நிலையானது என்பதால், வரையறுக்கும்
சமன்பாடுகளின் இடப்பக்கத்திலும் வலப்பக்கத்திலும் அது இடம்பெறும்.
\begin{align*}
\Add(x,0) & =  x\\
\Add(x,y+1) & =  \Add(x,y)+1
\end{align*}
இச்சமன்பாடுகள் அனைத்து $x$ \emph{மற்றும்}~$y$ க்குமான
$\Add$ இன் மதிப்பைக் குறிப்பிடுகின்றன. எடுத்துக்காட்டாக,
$\Add(2,3)$ ஐக் காண, $x = 2$ க்கான வரையறுக்கும் சமன்பாடுகளைப்
பயன்படுத்துகிறோம்: முதலாவதைக் கொண்டு $\Add(2,0) = 2$ ஐக் கண்டு,
பின்னர் இரண்டாவதை அடுத்தடுத்துப் பயன்படுத்தி
$\Add(2,1) = 2 + 1 = 3$, $\Add(2, 2) = 3 + 1 = 4$, $\Add(2,
3) = 4 + 1 = 5$ ஆகியவற்றைக் காண்கிறோம்.

$\Add$ இன் வரையறையில் இரண்டாம் சமன்பாட்டின் வலப்பக்கத்தில்
$+$ ஐப் பயன்படுத்தினோம்; ஆனால்~$1$ ஐக் கூட்ட மட்டுமே பயன்படுத்தினோம்.
வேறுவிதமாகக் கூறினால், அடுத்தெண் சார்பு $\Succ(z) = z+1$ ஐப்
பயன்படுத்தி, $\Add(x,y+1)$ ஐ வரையறுக்க முந்தைய மதிப்பு
$\Add(x,y)$ க்கு அதனைச் செயல்படுத்தினோம். எனவே முந்தைய மதிப்பின்மீது
நாம் செயல்படுத்தும் தனிச் சார்பைக் கொண்டு மீள்வரையறை தரப்பட்டுள்ளது
என்று கருதலாம். எனினும், அச்சார்பு முந்தைய மதிப்பை மட்டுமல்லாமல்
$x$ மற்றும்~$y$ ஐயும் சார்ந்திருக்க அனுமதிப்பதில் தீங்கில்லை;
சிலவேளைகளில் அது தேவைப்படுகிறது. பின்வருவதைப் பார்க்க:
\begin{align*}
  \Mult(x,0) & =  0 \\
  \Mult(x,y+1) & =  \Add(\Mult(x,y),x)
\end{align*}
முந்தைய மதிப்பு $\Mult(x,y)$ மற்றும் முதல் உள்ளீடு~$x$ ஆகிய
இரண்டிற்கும் சார்பு $\Add$ ஐச் செயல்படுத்துவதன் மூலம்
சார்பு $\Mult$ ஐ வரையறுக்கும் முதன்மை மீள்வரையறை இது.
அனைத்து உள்ளீடுகள் $x$ மற்றும்~$y$ க்குமான
சார்பு~$\Mult(x,y)$ ஐயும் இது வரையறுக்கிறது. எடுத்துக்காட்டாக,
$\Mult(2,0)$, $\Mult(2,1)$, $\Mult(2,2)$ மற்றும்~$\Mult(2,3)$
ஆகியவற்றை அடுத்தடுத்துக் கணிப்பதால் $\Mult(2,3)$ தீர்மானிக்கப்படுகிறது:
\begin{align*}
  \Mult(2,0) & = 0\\
  \Mult(2,1) & = \Mult(2,0+1) =
  \Add(\Mult(2,0), 2) = \Add(0, 2) = 2\\
  \Mult(2,2) & = \Mult(2,1+1) =
  \Add(\Mult(2,1), 2) = \Add(2, 2) = 4\\
  \Mult(2,3) & = \Mult(2,2+1) =
  \Add(\Mult(2,2), 2) = \Add(4, 2) = 6
\end{align*}

பொது முறை பின்வருமாறு: சார்பு~$h(x_0, \dots, x_{k-1}, y)$
ஒன்றின் முதன்மை மீள்வரையறையைத் தர, இரு சமன்பாடுகளை வழங்குகிறோம்.
முதலாவது,~$h$ ஐச் சுட்டாமல் $h(x_0, \dots, x_{k-1}, 0)$ இன்
மதிப்பை வரையறுக்கிறது. இரண்டாவது, $h(x_0,
\dots, x_{k-1}, y+1)$ இன் மதிப்பை $h(x_0, \dots, x_{k-1}, y)$,
பிற உள்ளீடுகள் $x_0$, \dots,~$x_{k-1}$ மற்றும்~$y$ ஆகியவற்றைக்
கொண்டு வரையறுக்கிறது. இரண்டாம் சமன்பாட்டில்~$h$ இன் உடனடியாக
முந்திய மதிப்பை மட்டுமே பயன்படுத்தலாம். இவ்விரு சமன்பாடுகளின்
வலப்பக்கங்கள் தரும் செயல்கள் தாமே சார்புகள் $f$ மற்றும்~$g$
என்று கருதினால், முதன்மை மீள்வரையறையால் புதிய சார்பு~$h$ ஒன்றை
வரையறுக்கும் பொது முறை இதுவாகும்:
\begin{align*}
  h(x_0, \dots, x_{k-1}, 0) & = f(x_0, \dots, x_{k-1})\\
  h(x_0, \dots, x_{k-1}, y+1) & =
  g(x_0, \dots, x_{k-1}, y, h(x_0, \dots, x_{k-1}, y))
\end{align*}  
$\Add$ இன் நேர்வில், $k=1$ மற்றும் $f(x_0) = x_0$
(ஒன்றாதல் சார்பு); மேலும் $g(x_0, y, z) = z + 1$ (தன் மூன்றாம்
உள்ளீட்டின் அடுத்தெண்ணைத் தரும் $3$-இடச் சார்பு):
\begin{align*}
  \Add(x_0, 0) & = f(x_0) = x_0\\
  \Add(x_0, y+1) & = g(x_0, y, \Add(x_0, y)) =
  \Succ(\Add(x_0, y))
\end{align*}
$\Mult$ இன் நேர்வில், $f(x_0) = 0$ (எப்போதும்~$0$ ஐத் தரும்
மாறிலிச் சார்பு) மற்றும் $g(x_0, y, z) = \Add(z,x_0)$
(தன் கடைசி மற்றும் முதல் உள்ளீடுகளின் கூட்டுத்தொகையைத் தரும்
$3$-இடச் சார்பு):
\begin{align*}
  \Mult(x_0,0) & =  f(x_0) = 0 \\
  \Mult(x_0,y+1) & =  g(x_0, y, \Mult(x_0,y)) =
  \Add(\Mult(x_0,y), x_0)
\end{align*}

\end{document}
